\section{Những Thư Viện, Công Cụ Hiện Có}

Hệ thống được xây dựng dựa trên các thư viện và công cụ mã nguồn mở phổ biến trong lĩnh vực xử lý tiếng nói và truy xuất thông tin. Các thành phần này có thể được sử dụng độc lập dưới dạng mô-đun hoặc kết hợp thành pipeline hoàn chỉnh:
\begin{itemize}[label={}]
    \item \textbf{Thư viện xử lý tiếng nói}: Các thư viện hỗ trợ phát hiện hoạt động giọng nói, trích xuất embedding và nhận dạng tiếng nói, cho phép xây dựng pipeline diarization và ASR theo từng bước.
    \item \textbf{Mô hình end-to-end}: Một số mô hình cung cấp khả năng xử lý trọn gói (ví dụ ASR end-to-end), được sử dụng như các khối chức năng thay thế nhằm so sánh và đánh giá tính linh hoạt của hệ thống.
    \item \textbf{Thư viện học máy và học sâu}: Các framework như PyTorch được sử dụng để tải và suy luận mô hình, đồng thời hỗ trợ thao tác tensor và tăng tốc bằng GPU.
    \item \textbf{Công cụ truy xuất và lưu trữ vector}: Các thư viện phục vụ việc xây dựng chỉ mục vector và truy vấn ngữ nghĩa, đóng vai trò trung tâm trong mô-đun RAG.
    \item \textbf{API mô hình ngôn ngữ lớn}: LLM được tích hợp thông qua API, cho phép thực hiện các tác vụ tổng hợp, trả lời câu hỏi và tái cấu trúc thông tin dựa trên ngữ cảnh truy xuất.
\end{itemize}